\documentclass[aspectratio=169, 10pt]{beamer}
\usepackage[utf8]{inputenc}
\usepackage[T1]{fontenc}
\usepackage{lmodern}
\usepackage{amsmath, amssymb}
\usepackage{circuitikz}
\usetheme{default}
\usecolortheme{dove}
\setbeamertemplate{navigation symbols}{}

\begin{document}
\begin{frame}{Funzione di trasferimento complessiva ($C_{\text{in}}$, $C_f$)}
\centering
\vspace{-0.2em}
\begin{circuitikz}[scale=0.78, transform shape, line width=0.7pt]
  % Ingresso Vin
  \draw (-1.2, 0.5) node[left] {$v_{in}(t)$} to[short, o-] (-0.2, 0.5);
  \draw (-0.2, 0.5) -- (-0.2, 1.6);
  \draw (-0.2, 0.5) -- (-0.2, -0.6);
  
  % Ramo parassita diretto Cp (in alto)
  \draw (-0.2, 1.6) to[C, l_={$C_p$}] (3.6, 1.6) -- (4.0, 1.6);

  % Ramo compensazione con blocco [-1] e Cc (in basso)
  \node[draw, rectangle, fill=white, inner sep=3.8pt, font=\scriptsize] (inv) at (1.0, -0.6) {$-1$};
  \draw (-0.2, -0.6) -- (inv.west);
  \draw (inv.east) -- (1.7, -0.6) to[C, l={$C_c$}] (3.6, -0.6) -- (4.0, -0.6);

  % Connessione al nodo di sense
  \draw (4.0, 1.6) -- (4.0, -0.6);
  \draw (4.0, 0.5) to[short, -*] (5.0, 0.5) coordinate (sense);
  
  % Cin verso massa
  \draw (5.0, 0.5) to[C, l={$C_{\text{in}}$}, *-] (5.0, -1.1) node[ground]{};

  % TIA OPA656 (distanziato da Cin per massima chiarezza visiva)
  \node[op amp] (opamp1) at (7.2, 0) {};
  \node[font=\scriptsize, below=3pt] at (opamp1.south) {OPA656};
  \draw (sense) -- (opamp1.-);
  \draw (opamp1.+) -- ++(0, -0.4) node[ground]{};

  % Retroazione TIA (Cf sopra Rf)
  \draw (sense) -- (5.0, 2.3) -- (5.7, 2.3);
  \draw (5.7, 3.0) to[C, l=$C_f$] (7.9, 3.0);
  \draw (5.7, 1.8) to[R, l=$R_f$] (7.9, 1.8);
  \draw (5.7, 3.0) -- (5.7, 1.8);
  \draw (7.9, 3.0) -- (7.9, 1.8);
  \draw (7.9, 2.3) -| (opamp1.out);

  % Secondo stadio AD817 (invertente G = -10)
  \node[op amp] (opamp2) at (11.5, -0.5) {};
  \node[font=\scriptsize, below=3pt] at (opamp2.south) {AD817};
  \draw (opamp1.out) -- (8.4, 0) to[R, l_={$R_1 = 100\,\Omega$}] (10.2, 0) -- (opamp2.-);
  \draw (opamp2.+) -- ++(-0.2,0) node[ground]{};

  % Retroazione AD817 (R2)
  \draw (opamp2.-) to[short, *-] ++(0, 1.5) to[R, l={$R_2 = 1\,\text{k}\Omega$}] ++(2.0, 0) -| (opamp2.out);
  
  % Uscita finale
  \draw (opamp2.out) to[short, -o] ++(0.8, 0) node[right] {$V_{\text{out}}$};
\end{circuitikz}

\vspace{0.8em}
{\small
Modellando l'operazionale OPA656 con $A(s) \approx \frac{\omega_t}{s} \implies V^- = -\frac{s}{\omega_t} V_{\text{out1}}$ \quad $\left(\text{e secondo stadio } V_{\text{out}} = -10\, V_{\text{out1}}\right)$
}

\vspace{0.4em}
\[
  \frac{V_{\text{out}}(s)}{V_{in}(s)} = \frac{10 \cdot s R_f (C_p - C_c)}{1 + s \left( R_f C_f + \frac{1}{\omega_t} \right) + s^2 \frac{R_f (C_{\text{in}} + C_p + C_c + C_f)}{\omega_t}}
\]
\end{frame}
\end{document}
